% Chapter 25: Light Is an Electromagnetic Wave
\chapter{Light Is an Electromagnetic Wave}

\step{A Counterintuitive Opening}

\begin{mainline}
Close your eyes and face the Sun for a while; even your eyelids feel warm. This is so ordinary that we never ask how it happens. Between you and the Sun lie a hundred and fifty million kilometers of emptiness: real emptiness, without air, water, rope, or even much dust. Yet something leaves the Sun, crosses that empty space for over eight minutes, reaches your skin, and warms it.

Think of more examples. Your indoor radio receives a station hundreds of kilometers away. WiFi passes through walls. Hospital X-rays pass through your body and leave bone shadows on film. Press a remote and the television changes channels. A firefighter's thermal camera sees a trapped person's body heat through thick smoke. These signals cross vacuum, walls, and bodies. What are they? Where there is nothing, what allows them to travel?

A subtler mystery hides in the living room. Flip a switch and the light comes on immediately; make a call and your voice immediately reaches someone a thousand kilometers away. Are these immediacies truly instantaneous? If so, what travels infinitely fast? If not, how fast does it travel, and what is traveling?

This chapter gives an unexpected answer, repeatedly confirmed by countless experiments: light, radio waves, and X-rays are all the same thing, a combined electric-and-magnetic disturbance sustaining itself as it races through space. What warms you in sunlight and what your radio receives belong to one family.
\end{mainline}

\step{Make a Guess First}

\begin{mainline}
Before reading on, write down three guesses.

First, does light take time to travel from Sun to Earth? If so, how long? Everyday experience says a room lights up as soon as you switch on a lamp, suggesting instantaneous arrival.

Second, water waves involve vibrating water, and sound involves vibrating air (Chapter 17). If light is a wave traveling through vacuum, what vibrates?

Third, electricity, magnetism, and light are old acquaintances met through frictional charging, compasses, and sunshine. Are they related?

Historically, almost everyone guessed that light arrived instantly. Galileo devised an honest experiment: two people with covered lanterns stood on hilltops one or two kilometers apart. A uncovered a lantern; on seeing it, B uncovered the other. A timed the interval until the reply arrived. Naturally, nothing measurable emerged, not because light takes no time, but because crossing a kilometer or two takes only a few millionths of a second, far beyond human reaction speed.

The real breakthrough came from the sky. In 1676, Danish astronomer Rømer noticed something odd: as Earth moved away from Jupiter, a Jovian moon's entries into Jupiter's shadow ran increasingly late, accumulating differences of twenty-odd minutes. His bold explanation was that the moon kept time, but light needed extra time to cross the growing distance. Its speed was finite, simply absurdly large and imperceptible over everyday distances. Increasingly precise measurements confirmed the guess: in one second, light could circle Earth's equator seven and a half times.

Keep the second and third guesses in mind. For the second, perhaps an invisible substance in vacuum vibrates, or perhaps some waves need no substance to vibrate at all. The latter sounds contradictory but stars in this chapter. For the third, history supplied an enticing clue: light speed could be calculated from constants measured in two purely electromagnetic experiments. If electricity, magnetism, and light were unrelated, there would be no reason for this coincidence. Now let us make some invisible light ourselves.
\end{mainline}

\step{Hands-on Experiment}

\begin{experiment}[Send a Telegraph Message with One Battery]
Materials: an AM radio, an AA battery, and a short piece of thick wire.

% BEGIN TRANSLATOR SAFETY NOTE
\textbf{Translator safety note:} Do not perform the direct battery-shorting or sparking procedure below, including its foil-wrapped repeat. Metal across the terminals can cause hazardous heating or leakage. Use a properly current-limited educational demonstration instead. See \href{https://oem.energizer.com/english/dos-donts/}{Energizer battery-care guidance}.
% END TRANSLATOR SAFETY NOTE

Tune between stations so the speaker produces only static. Hold one wire end against the battery's negative terminal and rapidly scrape the other across the positive terminal, making and breaking contact. You immediately hear clicks on the radio, even several meters away with no connecting wire.

What happened? At connection and disconnection, current changes abruptly, and the fields near the spark change with it. A small packet of electromagnetic waves is flung across the room, nudging electrons in the radio antenna. Its circuit amplifies that tiny push into a click.

You have repeated one of science's most famous experiments. In 1887, Hertz used larger sparks and more careful measurements to produce and receive electromagnetic waves, confirming Maxwell's paper prediction from more than twenty years earlier. Your tabletop setup is a miniature household version.

One extra step: wrap the entire radio in aluminum foil and scrape the wire again. The clicks disappear. The foil has not sucked anything away; the reason appears later.

This experiment contains all our main characters: changing current, the source; something crossing the room, the wave; and pushed electrons, reception. We must explain the middle link: by what law does that something cross the room without air or wires?
\end{experiment}

\begin{experiment}[Darkness between Two Polarizers]
Find two pairs of polarized sunglasses, or salvage two polarizing films from a discarded LCD. Stack them facing a lamp and slowly rotate one: transmitted light dims until it is nearly black at ninety degrees.

Now try something odd: insert a third polarizer at forty-five degrees between the two already-dark ones. The darkness brightens again!

If polarizers were simply light blockers, adding another could never increase light. Record the observation. Later it will help reveal a key conclusion: light's vibration is perpendicular to its direction of travel.
\end{experiment}

\step{The Old Model Runs into Trouble}

\begin{mainline}
Return to the official physical picture of the mid-nineteenth century.

Light was a wave, with strong evidence: it interferes and diffracts, as Chapter 35 explains, and only waves do that. Newton had favored light particles, neatly explaining straight-line propagation and reflection. But Young's alternating bright and dark interference fringes defeated that picture: adding two particle streams should give more particles, not darkness. Wave theory won, immediately inheriting a burden: what vibrates? Water moves in water waves, air in sound, but what moves when light crosses vacuum? Physicists invented the ether, filling the universe yet extraordinarily tenuous so planets could move through it for billions of years without drag. It also needed extraordinary rigidity to support light's rapid vibrations. A substance both imperceptibly tenuous and unimaginably rigid sounded increasingly implausible under scrutiny.

Meanwhile electricity and magnetism fought separately. Charges attracted and repelled across space; magnetic needles deflected across space. The prevailing account was action at a distance: forces arrived instantly, with no intervening process. It seemed sufficient for electrostatics, since Coulomb's law gave the right force; who cared how it arrived? But induction (Chapter 23) sounded an alarm. Insert a magnet into a coil and current immediately appears, as if something had arrived ahead of it. Faraday rejected action at a distance, imagining invisible field lines filling space, transmitting force and carrying changes in finite time. Most colleagues dismissed this as pictures from someone who lacked mathematics. Maxwell recognized their physical substance, translated them mathematically, and supplied the crucial displacement-current addition: a changing electric field substitutes for current in producing magnetic circulation (Chapter 24).

Then the equations spoke for themselves, saying something nobody expected: even in vacuum entirely devoid of charges and currents, a field can travel on its own.
\end{mainline}

\step{Inventing New Concepts}

\begin{mainline}
Maxwell asked what electric and magnetic fields his four equations allowed in vacuum.

Recall their meanings (Chapter 24): charge sources electric fields; isolated magnetic charges have never been found; changing magnetic fields create electric circulation; and current, plus Maxwell's changing-electric-field addition, creates magnetic circulation. Near wires and charges, fields are tied to sources. Far from every source, the third and fourth equations grip each other: magnetic circulation allows electric fields to vary in time, and electric-field change allows magnetic fields to vary. Could the pair maintain itself without any source?

The astonishing answer was yes. Suppose an electric field points vertically and varies horizontally. Together the third and fourth equations require it to vary with time, accompanied by a magnetic field perpendicular to both it and the propagation direction. The magnetic distribution also shifts forward as a whole. The complete solution resembles two perpendicular waving walls, electric and magnetic, moving hand in hand at a fixed speed. No wire, air, or ether is needed; vacuum is the whole stage.

Avoid a popular mistake: electric field first creates magnetic field, which creates electric field in a repeating forward relay. That is not what happens. The fields do not give birth to each other in sequence; they coexist in one solution and remain in step. Their joint spatial pattern moves forward with time, constrained throughout by Maxwell's equations. It resembles a stadium wave: each spectator stands and sits locally, while the wave races around the stands. The similarity is that a change of state propagates, not matter. The difference is that a stadium wave needs rows of people, while no people stand behind an electromagnetic wave. The field itself changes.

\begin{center}
\begin{tikzpicture}[>=Stealth,scale=0.85]
  \draw[->] (-0.2,0) -- (10.6,0) node[right]{Propagation};
  \draw[->] (0,-1.7) -- (0,2.0) node[above]{Electric field $E$};
  \draw[blue!70!black,thick,domain=0:10,samples=120] plot (\x,{1.25*sin(108*\x)});
  \foreach \x in {0.4,1.3,2.2,3.1,4.0,4.9,5.8,6.7,7.6,8.5,9.4}
    \draw[->,blue!50] (\x,0) -- (\x,{1.25*sin(108*\x)});
  \foreach \x in {0.83,4.17,7.5}
    \node[red!70!black] at (\x,-0.42) {\small$\odot$};
  \foreach \x in {2.5,5.83,9.17}
    \node[red!70!black] at (\x,-0.42) {\small$\otimes$};
  \node[red!70!black,anchor=west] at (7.6,-1.35) {\small Field $B$ ($\odot$ out, $\otimes$ into page)};
\end{tikzpicture}
\end{center}

More astonishing still, Maxwell calculated the propagation speed and found that only two constants determined it, one from electrostatic experiments, the other from magnetostatic experiments. Their measured values gave about 310,000 kilometers per second. The best light-speed measurement then, Fizeau's rotating-wheel measurement of 1849, was around 315,000 kilometers per second. Numbers from entirely different domains nearly coincided. Maxwell wrote in 1865 that such agreement gave ample reason to conclude that light itself was an electromagnetic disturbance.

Light, electricity, and magnetism, three old acquaintances, were one family.

The paper prediction still needed an experimental hand. In 1887, Hertz used an induction coil to spark between two copper balls, making current change violently in a tiny interval. Maxwell's equations said this should emit electromagnetic waves. Several meters away he placed a metal ring with a narrow gap. Sure enough, tiny but distinct sparks jumped across it. No wire, no contact: energy crossed the laboratory on its own. Hertz also reflected, interfered, and polarized the waves, checking their wave properties one by one, and measured their speed as light speed. Eight years after Maxwell's death, experiment vindicated his equations. A little over a decade later, Marconi turned similar sparks into transatlantic wireless telegraphy. Only one generation separated an equation from an industry.
\end{mainline}

\step{The Model in One Sentence}

\begin{oneline}
Light is a joint electric-and-magnetic disturbance, sustaining itself through vacuum at about 300,000 kilometers per second without any medium. Radio waves, microwaves, infrared, visible light, ultraviolet, X-rays, and gamma rays are all family members, differing only in how many times they oscillate each second.
\end{oneline}

\step{The Mathematical Translator}

\begin{mainline}
First see how speed emerges from the equations. Enter two constants: vacuum permittivity $\varepsilon_0\approx 8.85\times 10^{-12}$, from Coulomb's law, interpretable as how hard it is to establish an electric field in vacuum; and vacuum permeability $\mu_0 = 4\pi\times 10^{-7}\approx 1.26\times 10^{-6}$, from current's magnetic effects, describing its ability to establish a magnetic field. Maxwell's equations give
\[
  c=\frac{1}{\sqrt{\varepsilon_0\mu_0}}\,.
\]
Insert numbers: the product is $\varepsilon_0\mu_0\approx 1.11\times 10^{-17}$, its square root about $3.34\times 10^{-9}$, and its reciprocal
\[
  c\approx 3.0\times 10^{8}\ \text{m/s}\,.
\]
Constants measured entirely through electrostatic and magnetostatic experiments combine to equal light speed. This is one of physics' most beautiful collisions: two unrelated experimental domains meeting at one number.

Check immediately. If $\mu_0$ approaches zero, meaning a changing electric field no longer accompanies any magnetic field, the denominator vanishes and $c$ becomes infinite: waves collapse into instantaneous action at a distance. The same holds if $\varepsilon_0$ approaches zero. Neither constant, neither electrical nor magnetic hand, is dispensable. Conversely, $c$ cannot be zero because both constants are positive real numbers. For a directional check, reverse the whole electric field: the magnetic field also reverses, leaving propagation unchanged, just as expected for steadily forward energy transport.

Next comes every wave's indispensable relationship:
\[
  c=\lambda f\,,
\]
wavelength times frequency equals speed. Check extremes. As frequency $f$ approaches zero, wavelength $\lambda$ approaches infinity. A nearly time-independent field cannot propagate, returning to the static fields of Chapters 19 and 22, reasonably. At extremely high frequency, wavelength becomes short enough to enter atomic gaps, explaining X-rays imaging bones and gamma rays penetrating steel. Check everyday numbers too: 100-megahertz FM gives a wavelength of about 3 meters, comparable to your radio's telescoping antenna. WiFi at $2.4\times10^{9}$ hertz has wavelength about 12 centimeters, the same as the waves in a microwave oven: they occupy neighboring bands. Visible light with wavelength $\lambda\approx 500$ nanometers gives frequency $f = 3\times10^{8}/(5\times10^{-7}) = 6\times10^{14}$ hertz, about 600 trillion oscillations each second. No wonder nobody directly sees the electric field oscillate; we detect its average effects on matter, such as heating, photographic response, and antenna-electron motion.

Arrange frequency from low to high and obtain the electromagnetic family portrait. One equation, one speed, different numbers of oscillations per second:
\end{mainline}

\begin{center}
\begin{tikzpicture}[>=Stealth]
  \fill[blue!10]   (0,0)     rectangle (4.5,0.55);
  \fill[cyan!15]   (4.5,0)   rectangle (6.75,0.55);
  \fill[green!15]  (6.75,0)  rectangle (8.7,0.55);
  \fill[yellow!50] (8.7,0)   rectangle (8.925,0.55);
  \fill[orange!20] (8.925,0) rectangle (10.5,0.55);
  \fill[red!15]    (10.5,0)  rectangle (12,0.55);
  \fill[violet!18] (12,0)    rectangle (13.5,0.55);
  \draw (0,0) rectangle (13.5,0.55);
  \foreach \x in {4.5,6.75,8.7,8.925,10.5,12} \draw (\x,0) -- (\x,0.55);
  \node at (2.25,0.85) {\small Radio waves};
  \node at (5.6,0.85) {\small Microwaves};
  \node at (7.7,0.85) {\small Infrared};
  \node at (9.7,0.85) {\small Ultraviolet};
  \node at (11.25,0.85) {\small X-rays};
  \node at (12.75,1.05) {\small \shortstack{Gamma\\rays}};
  \draw[->] (8.81,-0.15) -- (8.81,-0.6);
  \node[anchor=north] at (8.81,-0.6) {\small Visible light: this narrow strip};
  \foreach \x/\l in {0/$10^{3}$, 4.5/$10^{9}$, 6.75/$10^{12}$, 10.5/$10^{17}$, 12/$10^{19}$, 13.5/$10^{21}$}
    \node[below] at (\x,0) {\scriptsize \l};
  \node[below] at (6.75,-1.4) {\scriptsize Frequency / Hz};
\end{tikzpicture}
\end{center}

\begin{mainline}
Meet the family. Radio waves have the lowest frequencies and longest wavelengths, from meters to kilometers, serving broadcasts, walkie-talkies, and satellite television. Microwaves are shorter: WiFi uses about $2.4\times10^{9}$ and $5\times10^{9}$ hertz, while ovens use $2.45\times10^{9}$ hertz, wavelengths of a dozen or so centimeters. Infrared is shorter still, encompassing remote-control LEDs, night-vision instruments, and heat radiation from your face. A little everyday magic: point a phone camera at a television remote's invisible infrared LED and press a button. The screen shows flashes because its silicon sensor responds to near infrared, while your eyes do not. Visible light occupies only the narrow slit, about 380 to 760 nanometers. Ultraviolet can break chemical bonds, tanning skin and killing bacteria; the blue-violet lamp in a sterilizer relies on it. X-rays pass through soft tissue but are blocked by bone, giving medical images. Gamma rays come from nuclear changes, with the highest frequencies and greatest penetration, used in radiotherapy and archaeological dating.

Neither end of the chart is a final boundary. Electromagnetic waves above gamma-ray frequencies can exist in principle, though they become harder to produce and detect; lower frequencies than radio can too. This is not seven different things standing in line, but a continuous ruler for one thing. We name its regions by what happens when they interact with us.

An engineering application: radar. A station emits a microwave pulse, receives its aircraft reflection, and multiplies round-trip time by $c$, dividing by two to obtain distance. Police speed detectors use the same idea, additionally examining a small echo-frequency shift. Navigation satellites reverse the arrangement: several continuously broadcast their positions and transmission times; your phone compares arrival times to infer its location on Earth. The foundation is electromagnetic waves traveling straight at constant speed $c$. Giant radio telescopes transmit nothing, merely raising large dishes to receive waves quietly from distant celestial objects. Much of humanity's knowledge of the universe comes from listening to this family's different members.
\end{mainline}

\begin{mathbridge}
Why exactly $1/\sqrt{\varepsilon_0\mu_0}$? A simplified one-dimensional model reveals the structure.

Let the electric field have only a $y$ component and vary only along $x$, written $E(x,t)$, while the magnetic field has only a $z$ component, $B(x,t)$. Place a narrow rectangular loop in the $x$-$y$ plane and apply Chapter 23's Faraday law: the electric-field difference between its long sides equals the magnetic-flux change rate through it. Shrink its width to zero:
\[
  \frac{\partial E}{\partial x}=-\frac{\partial B}{\partial t}\,.
\]
Now place the loop in the $x$-$z$ plane and apply Ampère's law with displacement current, in vacuum without conduction current. The same limit gives
\[
  -\frac{\partial B}{\partial x}=\varepsilon_0\mu_0\,\frac{\partial E}{\partial t}\,.
\]
Differentiate the first again with respect to $x$ and substitute the second, combining them:
\[
  \frac{\partial^2 E}{\partial x^2}=\varepsilon_0\mu_0\,\frac{\partial^2 E}{\partial t^2}\,.
\]
Compare Chapter 17's standard wave equation: second spatial derivative equals second time derivative divided by speed squared. Immediately,
\[
  c=\frac{1}{\sqrt{\varepsilon_0\mu_0}}\,.
\]
Repeat the same two steps for $B$ to obtain the identical equation and speed. The electric and magnetic waves are locked into one procession; neither can go faster or slower.
\end{mathbridge}

\begin{challenge}
Three dimensions require no new physics, only vector calculus. In source-free vacuum, take the curl of Faraday's law $\nabla\times\bm E=-\partial\bm B/\partial t$. Use $\nabla\times(\nabla\times\bm E)=\nabla(\nabla\cdot\bm E)-\nabla^2\bm E$ and the vacuum condition $\nabla\cdot\bm E=0$, then substitute the Ampère--Maxwell law:
\[
  \nabla^2\bm E=\varepsilon_0\mu_0\,\frac{\partial^2\bm E}{\partial t^2}\,,
\]
with the identical equation for the magnetic field. Each Cartesian component of $\bm E$ thus satisfies an independent three-dimensional wave equation.

Transversality follows directly too. For a plane wave $\bm E=\bm E_0\cos(\bm k\cdot\bm r-\omega t)$, the condition $\nabla\cdot\bm E=0$ gives $\bm k\cdot\bm E_0=0$: oscillation is perpendicular to propagation. The condition $\nabla\cdot\bm B=0$ gives the same for magnetism. Notice that no reference frame appears anywhere in the equation. Relative to whom is $c$ measured? Keep that question; it opens Chapter 38.
\end{challenge}

\step{The Model's Limits}

\begin{mainline}
The electromagnetic-wave model is astonishingly powerful, but bounded.

First: nearby fields are not waves. Close to antennas, transformers, and charging coils, field patterns are complex and decay rapidly with distance. Most energy circles back into the circuit instead of escaping. Only several wavelengths from a source do fields settle into freely propagating waves. Thus calling every changing field an electromagnetic wave goes too far; only the part genuinely flung away qualifies.

Second: speed changes in media. In glass, water, and air, electric fields continually shake material charges, reducing wave speed to $c/n$, where $n$ is the refractive index. Light in water travels at only about $0.75\,c$. The constant in Maxwell's equations is vacuum speed $c$; Chapter 36 covers media.

Third: very weak light breaks this model. Dim it repeatedly and detectors begin responding one event at a time, instead of receiving a continuously weakening wave. Smooth electromagnetic fields are an excellent model, not the final word. Light also has a particulate side, the quantum story beginning in Chapter 42. Distinguish levels: light with many photons behaving like a classical wave is an experimental fact; Maxwell's equations describing waves are a repeatedly tested mathematical model; light being a vibrating field is our most useful physical explanation. These are not the same statement.

Fourth: the unanswered question. Relative to what is $c=1/\sqrt{\varepsilon_0\mu_0}$? Water-wave speed is relative to water; sound speed is relative to air. Maxwell's equations contain no term specifying where ether is at rest. This is no minor gap: it leads directly to special relativity (Chapter 38).

Three typical misuses deserve mention. First, imagining waves as ripples dispersed by air: on the contrary, air complicates matters, and vacuum provides the cleanest propagation. Second, saying microwave ovens use water's resonant frequency: $2.45\times10^{9}$ hertz is no sharp water resonance, but an engineering compromise between heating efficiency, penetration depth, and frequency allocation. Third, panicking at the word radiation. This model identifies the distinction: frequency. Below visible light, an individual wave packet lacks energy to break chemical bonds and affects the body chiefly through heating. Ultraviolet and above can break bonds and damage molecules. Radiation is not one danger label, but different marks on a frequency ruler.
\end{mainline}

\begin{mainline}
Two more important conclusions follow directly.

First, antennas. A transmitting antenna is a wire accelerating charges back and forth. Uniformly moving charges merely carry their fields along; accelerated charges can fling fields away as free waves. Efficient emission requires antenna size comparable to wavelength, hence medium-wave broadcasting towers over a hundred meters tall for wavelengths around 300 meters, telescoping radio antennas tens of centimeters long for FM wavelengths around 3 meters, and phone antennas a few centimeters long hidden inside devices operating near $2\times10^{9}$ hertz with wavelengths around 15 centimeters. Reception involves no mysterious signal-sucking force: a passing electric field pushes free metal electrons back and forth, producing a tiny current that an amplifier reads. Transmission and reception are two directions of one law. Phone signal bars reflect how strongly the received field pushes antenna electrons. Distant base stations, thick walls, or a hand over the antenna weaken the push. Elevator reception fails because its metal cabin becomes a shielding enclosure, just like foil around your radio.

\begin{center}
\begin{tikzpicture}[>=Stealth]
  \draw[very thick] (-2.2,0) -- (-0.25,0);
  \draw[very thick] (0.25,0) -- (2.2,0);
  \fill (0,0) circle (1.6pt);
  \draw[<->] (-1.3,0.4) -- (1.3,0.4) node[midway,above]{\small Charges oscillate};
  \draw[decorate,decoration={snake,amplitude=1.5pt,segment length=7pt},->] (2.5,0) -- (4.9,0) node[right]{\small Electromagnetic wave};
  \node[below] at (0,-0.25) {\small Transmitting antenna};
\end{tikzpicture}
\end{center}

Second, polarization and transverse waves. Electric oscillation is perpendicular to propagation, but can point in different transverse directions: vertical, horizontal, or diagonal. This is polarization. Longitudinal waves cannot do this: sound oscillates along its propagation direction and looks the same under rotation around that axis, so there is no polarization. Light's ability to polarize is therefore strong evidence of transversality. Sunlight is unpolarized, mixing orientations. A polarizer resembles a row of fine fences, passing one field component. Cross two polarizers and the second blocks everything passed by the first, producing the experiment's darkness. The analogy's limit: a polarizer selects oscillation direction like a fence, but the rejected component does not squeeze between slats; the material absorbs it as heat. Polarized fishing sunglasses use this principle: water reflections are mainly horizontally polarized, so vertically transmitting lenses remove most glare and reveal fish below. Why does the diagonal third sheet revive some light? Each sheet projects the incident field onto its transmission direction, rather than simply reducing light. The forty-five-degree middle sheet takes a diagonal component from the first sheet's output; the last then extracts its own component from that diagonal field. After two projections, some remains. More blockers mean more darkness, but more projectors can mean more light. That is the intuition trap; Chapter 36 gives quantitative Malus's law.
\end{mainline}

\step{Explaining the Opening Phenomenon}

\begin{mainline}
Return to sunlight. In the Sun's core and surface, thermal motion continually accelerates charged particles, which continually fling electromagnetic waves into space. A small fraction heads toward Earth, crossing a hundred and fifty million kilometers of vacuum in about eight minutes twenty seconds. Nothing carries it along; it sustains itself and shifts forward according to Maxwell's equations. At your skin, its electric field shakes molecular charges, and their accumulated energy becomes heat, an account for Chapters 26 through 28. The warmth on your eyelids is electric and magnetic field that left the Sun eight minutes ago, now pushing charges in your face.

Radios, WiFi, remotes, and X-ray images are frequency variants of the same phenomenon. The experiment's click is clear too: the battery spark emits broadband waves containing frequencies the radio receives. Wrapping it in foil makes the incoming field rearrange free electrons in the foil, generating an opposing field inside that nearly cancels it. The wave is not absorbed away, but pushed back by the conductor's own field.

The opening mystery is answered: neither a particle stream nor ripples in a mysterious medium cross vacuum, but the field itself.

The supposedly instantaneous lamp is explained too. Light needs only nanoseconds to reach your eyes, too brief for any human sense. Immediate merely means faster than everyday experience can resolve. Across astronomical distances, finite speed becomes visible fact: the Moon we see is over a second old, the Sun eight minutes old, and Andromeda two and a half million years old. Looking at the stars is watching a historical documentary edited by light speed, another starting point for Chapter 38's reconsideration of space and time.
\end{mainline}

\step{Thought Experiments and Challenges}

\begin{thought}[Put a Lid on the Sun]
If the Sun suddenly went out now, would we know immediately? No. The last sunlight would still travel for eight minutes twenty seconds; Earth would carry on normally, with no earlier message possible. More remarkably, even a gravitational change announcing the Sun's disappearance could not arrive faster: gravity's changes also propagate at light speed, a preview of Chapter 41.

Push the scale: call Proxima Centauri, 4.2 light-years away, and the interval between saying hello and hearing hello back exceeds eight years. The Milky Way is about a hundred thousand light-years across; a galaxy-wide video conference delays every sentence by a hundred thousand years. Electromagnetic waves are the universe's fastest messengers, yet the universe makes even them look painfully slow.
\end{thought}

First an estimate using real numbers, then three challenges.

% BEGIN TRANSLATOR SAFETY NOTE
\textbf{Translator safety note:} Do not disable, obstruct, or improvise changes to a microwave's turntable, controls, shielding, or interlocks. Follow the manufacturer's operating instructions; otherwise keep this calculation theoretical. Hot food and containers can burn. See \href{https://www.fda.gov/radiation-emitting-products/resources-you-radiation-emitting-products/microwave-ovens}{FDA microwave guidance}.
% END TRANSLATOR SAFETY NOTE

\textbf{Estimate: measure light speed with a tray of cheese.} Stop the microwave turntable, or support the food so it stays still. Lay out cheese slices or a chocolate bar and heat at low power for a dozen seconds or so. Melted hot spots appear, separated by cooler areas. These are cavity standing-wave antinodes, with adjacent spots half a wavelength apart (Chapter 17). Measure about 6 centimeters, giving wavelength $\lambda\approx 12$ centimeters. The appliance label specifies $f=2450$ megahertz $=2.45\times10^{9}$ hertz. Thus
\[
  c=\lambda f\approx 0.12\times 2.45\times10^{9}\approx 2.9\times10^{8}\ \text{m/s}\,,
\]
only about $3\%$ from the accepted value. A tray of cheese measures a fundamental universal constant, direct everyday evidence that light is an electromagnetic wave.

\begin{puzzles}
\item A microwave's glass door contains metal mesh with holes only a millimeter or two wide. You see the food clearly, but microwaves cannot escape. Why? Would replacing the mesh by a slit 20 centimeters long and 1 millimeter wide still be safe?\\
Hint: passage through a hole depends on wavelength relative to its largest dimension, not on whether the material looks transparent. What happens when the slit's long side exceeds the microwave wavelength?

% BEGIN TRANSLATOR SAFETY NOTE
\textbf{Translator safety note:} The mesh-change question is theoretical. Never cut, replace, or modify microwave shielding, door seals, or interlocks; do not operate an oven with a damaged door. See \href{https://www.fda.gov/radiation-emitting-products/resources-you-radiation-emitting-products/microwave-ovens}{FDA microwave safety guidance}.
% END TRANSLATOR SAFETY NOTE

\item Put a phone during a call inside a tightly closed metal lunchbox and reception stops; a plastic box leaves it unchanged. Yet a phone inside an old microwave oven that is not powered sometimes still receives calls. Explain.\\
Hint: shielding comes from induced conductor charges and currents generating an opposing field. Gaps comparable to wavelength or poor contacts let this cancellation system leak. Phone wavelengths range from a dozen to several dozen centimeters, and an old oven's door seal may not close perfectly everywhere.

% BEGIN TRANSLATOR SAFETY NOTE
\textbf{Translator safety note:} Do not put a phone into a microwave for a home test; never power an oven with electronics inside. Phone reception is not a validated microwave-leakage test. Consult the manufacturer about suspected damage or leakage. See \href{https://www.fda.gov/radiation-emitting-products/resources-you-radiation-emitting-products/microwave-ovens}{FDA microwave safety guidance}.
% END TRANSLATOR SAFETY NOTE

\item In a long road tunnel, FM with wavelength about 3 meters often retains some reception while AM with wavelength about 300 meters disappears first. But long waves supposedly bend around mountains and buildings better. Is that contradictory?\\
Hint: regard the entrance as an aperture. Much larger than a wavelength, it allows easy passage; much smaller, it admits little. Diffraction goes around obstacles; a tunnel requires entering an opening. These are different questions. This is open-ended: organize your reasoning by comparing wavelength with obstacle dimensions.

\item On an old radio receiving AM, rotating the telescoping antenna noticeably changes volume, whereas FM is less affected. Guess why, and consider why lengthening an antenna sometimes only seems to improve reception.\\
Hint: the field pushes electrons along its direction. An antenna aligned with it receives the strongest push; perpendicular alignment receives almost none. Distant AM stations have comparatively fixed field direction, whereas urban reflections scramble FM arrival directions. Lengthening really changes how antenna dimensions match wavelength.
\end{puzzles}
